\chapter{Review of Literature}
\section{Real-Time Audio-to-Text Systems}
Real-time audio-to-text systems are designed to convert spoken language to written text with minimal latency, a critical capability, and diverse applications in education, healthcare, and governance \cite{Kheddar2024AdvancedDL}. Recent advances in Automatic Speech Recognition (ASR) technology, largely driven by deep learning models and end-to-end (E2E) architectures such as transformers, have significantly improved the accuracy and reliability of these systems \cite{Sainz2024Realtime}. The ability to transcribe speech instantaneously is vital for applications such as live media captioning and interactive voice assistants, which require that systems maintain a Real-Time Factor (RTF) less than 1 \cite{Sainz2024Realtime}. State-of-the-art ASR models have achieved near-human performance in ideal audio conditions, typically demonstrating Word Error Rates (WER) below $5\%$ on common English benchmarks like LibriSpeech clean data \cite{Radford2023Whisper}. A key development in model performance and language coverage is Whisper (trained on 680,000 hours of multilingual data), which demonstrated improved robustness and high accuracy across over 90 languages. However, performance often remains variable for low-resource languages (e.g. Pashto, Punjabi) and in noisy environments, prompting ongoing research into fine-tuning and domain adaptation to maintain high accuracy universally \cite{Sehar2025LowResource}.

\noindent However, the inherent multilingual nature of environments such as Kenya, characterized by the prevalence of code switching and translanguaging among English, Kiswahili, and various indigenous languages, presents unique challenges to ASR systems. These complexities necessitate the development of high-quality datasets that comprehensively capture instances of code-switching and mixed-language speech, which are notably scarce in East Africa. Recent research emphasizes the critical need for robust ASR models capable of processing languages with diverse phonetic characteristics and limited digital resources \cite{s23229096}.

\noindent Studies indicate that deep learning-based approaches have established that multilingual Automatic Speech Recognition (ASR), particularly through transfer learning, offers a proven path to developing systems for languages with limited data. For instance, in their comprehensive survey of low-resource speech recognition, Yu et al. (2020) confirmed that models pre-trained on high-resource languages (HRLs) and fine-tuned on low-resource languages (LRLs) consistently achieved significant Word Error Rate (WER) reductions. However, their overview also critically highlighted two persistent limitations: the difficulty of knowledge transfer due to typological mismatch between source and target languages, and the need for the field to establish real-world standardized LRL benchmarks instead of relying on simulated data scarcity \cite{9180290}. Similarly, Mirishkar et al. (2021) investigated hybrid architectures for low-resource multilingual speech recognition in the Indian context, demonstrating approaches that could be relevant to other low-resource settings. The importance of large-scale multilingual corpora for training these models is consistently highlighted, as fine-tuning on language-specific datasets can lead to high accuracy in recognizing speech across languages like English and Kiswahili. Dong et al. (2023) proposed a speech recognition method based on domain-specific datasets and confidence decision networks, which could contribute to improving ASR performance in specialized contexts.

\noindent Beyond mere transcription, the integration of Natural Language Processing (NLP) capabilities, such as summarization, paraphrasing, and sentiment analysis, is crucial for achieving deeper linguistic understanding.Recent trends in deploying real-time ASR systems increasingly emphasize lightweight, API-driven architectures to enable scalable and low-latency access to speech services across platforms.


\section{Sentiment Analysis and Emotion Detection}
Sentiment analysis plays a pivotal role in understanding the emotional tone conveyed through spoken or written content, encompassing the identification and categorization of emotions expressed in text or speech. As a prominent area within NLP, it effectively mines implicit emotions in text, thereby assisting enterprises and organizations in making informed decisions \cite{RahmanJim2024}. Integrating sentiment analysis with audio-to-text systems will allow for a more nuanced understanding of communication, providing insights into user engagement and emotional responses, aligning with recent trends in deep learning-based NLP \cite{Sajid2023}.

\noindent Recent advances in machine learning and NLP have facilitated the development of sophisticated sentiment analysis models. Transfer learning has emerged as a key technique that utilizes existing knowledge to solve problems across different domains and produces state-of-the-art prediction results in sentiment analysis. The current academic landscape for sentiment analysis (SA) is heavily focused on the textual modality, demonstrating both advanced capabilities and persistent limitations. Comprehensive studies, such as the exploration by \cite{chinedu2023unraveling}, review the spectrum of contemporary SA techniques, ranging from foundational approaches to sophisticated deep learning models designed for unraveling emotions from written content.

\noindent This research has found profound application across various high-impact domains. For instance, in the demanding realm of financial markets, \cite{Jawale_Jawale_Ingale_Shetty_2023} successfully demonstrated the utility of SA for prediction and decision support. However, their work also underscored common methodological challenges, particularly the substantial requirement for large, high-quality datasets and the need for robust strategies to mitigate inherent model biases. Beyond finance, NLP techniques have proven vital in social contexts: \cite{9781308} highlighted their use in analyzing textual feedback in educational environments to identify areas for improvement through sentiment annotations, topic modeling, and summarization. Furthermore, the application of NLP extends to extracting crucial mental health insights from complex, real-world social media data \cite{10.2196/68720}, illustrating its capacity for handling noisy, unstructured linguistic inputs.

\noindent Despite these functional strides, the reviewed literature reveals shared and significant hurdles. Shaik et al. (2022) highlighted the persistent complexity of handling context-based challenges such as sarcasm, domain-specific language, and ambiguity issues that often require intricate feature engineering. Compounding this challenge, the systematic review of multilingual sentiment analysis by \cite{Abdullah_Rusli_2021} noted that English remains the overwhelmingly dominant language in the field, exposing a critical resource and model gap for global applications. 

\noindent Critically, the uniform focus on the written word in all the reviewed literature confirms a significant limitation: the field's current state of the art largely ignores the non-verbal emotional cues present in spoken communication. By relying solely on transcribed text, these models inherently fail to capture the full emotional context delivered through acoustic features (such as tone, pitch, and speed). This highlights the limitations of text-only sentiment analysis and motivates research into more context-aware approaches, including improved linguistic modeling of transcribed speech within integrated audio-to-text systems.


\noindent In the context of Kenya and East Africa, where multilingualism is prevalent, the ability to analyze sentiments in both Kiswahili and English is of paramount importance. This capability can provide valuable insights into user engagement, emotional responses, and interpersonal dynamics, for example, in customer service applications, real-time sentiment analysis can help businesses gauge customer satisfaction, allowing for immediate intervention and issue resolution. The continuous growth of data presents both opportunities and challenges for sentiment analysis, particularly in diverse linguistic environments.

\section{Text Summarization in Low-Resource and Speech-Based NLP Systems}

Text summarization is a fundamental Natural Language Processing (NLP) task aimed at condensing large volumes of text into concise representations while preserving key semantic information. In the context of speech-based systems, summarization plays a critical role in transforming lengthy spoken content such as lectures, meetings, or civic forums into digestible and actionable insights. This capability is particularly valuable in low-resource and multilingual environments, where access to information may be constrained by language barriers, literacy levels, or time limitations.

Modern text summarization approaches are broadly categorized into extractive and abstractive methods. Extractive techniques identify and select salient sentences directly from the source text, while abstractive approaches generate novel sentences that paraphrase and compress the original content. Recent advances in deep learning have shifted the research focus toward abstractive summarization, primarily driven by transformer-based architectures that model long-range dependencies and semantic context effectively \cite{raffel2020exploring}.

Among these architectures, the Text-to-Text Transfer Transformer (T5) has emerged as a state-of-the-art framework by unifying all NLP tasks including summarization into a single text-to-text format. Studies have demonstrated that T5 achieves strong performance across multiple summarization benchmarks by leveraging large-scale pretraining followed by task-specific fine-tuning \cite{raffel2020exploring}. Multilingual extensions of such models have further enabled cross-lingual transfer, allowing knowledge learned from high-resource languages to benefit low-resource languages such as Kiswahili.

However, summarization in low-resource settings introduces unique challenges. The scarcity of large, high-quality annotated datasets limits direct supervised training, necessitating reliance on transfer learning and multilingual pretraining. Additionally, summarization models applied to ASR-generated text must contend with transcription errors, disfluencies, and informal speech patterns, which can significantly degrade summary quality if not adequately addressed \cite{huang2021leveraging}. These challenges are amplified in real-time systems, where latency constraints restrict the use of computationally intensive post-processing techniques.

Evaluation of text summarization systems typically relies on automated metrics such as ROUGE and BLEU, which measure lexical overlap between generated summaries and reference texts. While these metrics have limitations particularly in capturing semantic equivalence they remain widely adopted due to their reproducibility and suitability for benchmarking abstractive summarization models \cite{lin2004rouge,papineni2002bleu}. Their use is especially practical in low-resource contexts where human evaluation may be costly or infeasible.

In multilingual African contexts, including Kenya, the application of summarization to transcribed Swahili speech offers significant potential for improving information accessibility in education, governance, and public discourse. Despite this promise, existing literature reveals a notable gap in research focusing on abstractive summarization for Swahili, particularly when applied downstream of Automatic Speech Recognition systems. This gap motivates the inclusion of a summarization component in integrated speech intelligence pipelines tailored to low-resource languages.


\section{Challenges in Multilingual Speech Recognition}
Despite significant advancements, several challenges persist in sentiment analysis technologies, particularly in multilingual settings. A primary obstacle is the scarcity of high-quality annotated audio-text datasets that accurately reflect the diversity of speech patterns and dialects across different regions. In East Africa, for example, the lack of publicly available datasets for languages such as Kiswahili poses a substantial barrier to developing robust sentiment analysis and ASR systems.

\noindent Another prominent challenge in multilingual societies, such as Kenya, is code switching and translanguaging. Code-switching complicates the development of ASR systems because models must be able to detect and process multiple languages within a single conversation. Current ASR systems often struggle with accurately transcribing code-switched speech, as they are typically trained on monolingual datasets \cite{s23229096}. To address this, researchers have explored approaches such as incorporating language identification modules into ASR systems, enabling dynamic language switching based on detected speech patterns.

\noindent Furthermore, cultural and contextual nuances in language present additional challenges for sentiment detection systems. These systems must be trained to recognize the emotional context of specific words and phrases within a particular cultural framework, which is especially critical in diverse regions like East Africa. The limited availability of labeled sentiment data in languages such as Kiswahili further hinders the development of effective sentiment analysis systems tailored to local contexts. Research on sentiment analysis for code-mixed text underscores the scarcity of research and datasets for such scenarios in low-resource languages, highlighting the need for dedicated platforms to investigate these problems.

\section{Educational and Social Implications}
The integration of real-time audio-to-text and sentiment analysis technologies has profound implications for education and social communication in multilingual societies. In educational settings, these technologies can significantly improve language learning by providing immediate feedback on language skills and emotional engagement. Real-time transcriptions of lectures and discussions can improve accessibility for students with hearing impairments or those who speak different languages. Furthermore, sentiment analysis can provide information on how students emotionally respond to the learning environment, allowing educators to adjust their methods and improve student outcomes. Research into applying transfer learning to sentiment analysis in social networks has shown that this approach can lead to better performance with limited data, a crucial factor for developing effective educational tools where comprehensive datasets may be scarce.

\noindent In social communication, sentiment analysis can help bridge communication gaps in various contexts, from government services to healthcare. For example, in public health campaigns, AI-enabled analysis of public attitudes on platforms like Facebook and Twitter can ensure messages effectively reach diverse populations, allowing for better engagement and understanding. Hu et al. (2025) conducted a scoping review of natural language processing technologies for public health in Africa, highlighting the potential of such tools to address health communication challenges. In governance, these technologies can improve citizen participation by providing insights into public sentiment and concerns, as sentiment analysis is a crucial process to gather and analyze people's opinions to aid decision making for corporations, governments and individuals.

\noindent The integration of real-time audio-to-text and multilingual sentiment analysis technologies presents significant opportunities to improve communication, learning, and accessibility, particularly in multilingual contexts. Although challenges related to data set availability, code switching, and cultural nuances remain, advances in these technologies continue to push the boundaries of what is possible. As these technologies evolve, their potential to transform education, communication, and social interaction in East Africa and beyond becomes increasingly evident. The ongoing development of robust, multilingual solutions will be crucial in bridging linguistic divides, fostering inclusivity, and driving innovation across sectors.

\subsection{Scientific Gap}

Despite notable advancements in Automatic Speech Recognition (ASR), sentiment analysis, and text summarization for high-resource languages, and foundational progress for low-resource languages, a significant scientific gap persists. While individual components of ASR, sentiment classification, and abstractive summarization have each demonstrated promise in isolation, there is a distinct lack of research focusing on their integration into a unified, real-time audio-to-text intelligence pipeline specifically optimized for Swahili.

Existing solutions often suffer from one or more limitations: inadequate transcription accuracy for spontaneous or code-switched Swahili speech, sentiment models that rely solely on transcribed text without accounting for linguistic and cultural nuance, or summarization systems that are not robust to transcription noise generated by ASR models. Furthermore, most prior studies evaluate these components independently, without considering the compounding effects of error propagation across an end-to-end system.

The challenges of code-switching, translanguaging, and culturally grounded sentiment expression in Swahili remain largely underexplored within integrated real-time systems. This gap underscores the need for a holistic approach that jointly considers ASR accuracy, sentiment reliability, summarization quality, and system latency within a single deployable framework. Addressing this gap is essential to enable practical, real-world speech intelligence applications in multilingual African contexts.

% \subsection{Scientific Gap}
% Despite the notable advancements in Automatic Speech Recognition (ASR) and sentiment analysis for high-resource languages, and the foundational work initiated for low-resource languages, a significant scientific gap persists. While individual components (ASR, sentiment analysis) have shown promise, there is a distinct lack of research focusing on integrated, real-time audio-to-text sentiment analysis pipelines specifically designed and optimized for the unique linguistic complexities and data scarcity of Swahili. Existing solutions often either lack the necessary accuracy and robustness for practical, real-time deployment in a multilingual African context or do not comprehensively address the end-to-end process from raw audio to actionable sentiment insights. Furthermore, the challenges of code-switching, translanguaging, and culturally nuanced sentiment expression in Swahili remain largely underexplored within such integrated real-time systems. This gap highlights the urgent need for a holistic approach that considers both the technical integration of models and the specific linguistic and contextual characteristics of Swahili to enable effective communication analysis.

\subsection*{2.6. Lines of Research} Stemming from the identified scientific gap, several lines of research can be pursued to advance the field of low-resource language processing, particularly for Swahili: \begin{enumerate} \item \textbf{Development of Enhanced Data Collection and Annotation Frameworks:} One crucial avenue involves creating systematic and scalable methodologies for collecting and annotating large, high-quality datasets for Swahili audio and text, specifically focusing on diverse accents, dialects, and instances of code-switching. This would involve community-driven annotation efforts and leveraging existing linguistic resources. \item \textbf{Exploration of Novel Multilingual and Cross-Lingual Transfer Learning Architectures:} Further research could focus on designing and evaluating new deep learning architectures that are inherently more robust to multilingual inputs and can effectively transfer knowledge from high-resource languages to low-resource ones with minimal fine-tuning data. This might include exploring advanced self-supervised learning techniques. \item \textbf{Integration of Linguistic Features for Improved Contextual Understanding:} Another line of research could involve incorporating explicit Swahili linguistic features (e.g., morphological analysis, syntactic structures) directly into the neural network models to enhance their understanding of context and nuance, particularly for sentiment detection in complex sentences. \item \textbf{Development of Robust Evaluation Metrics for Low-Resource Contexts:} Given the unique challenges, there is a need to develop and validate new evaluation metrics or adapt existing ones to more accurately assess the performance of ASR and sentiment analysis systems in low-resource, code-switched environments. \item \textbf{Real-Time End-to-End Pipeline Optimization:} A critical research direction involves optimizing the entire audio-to-text sentiment analysis pipeline for real-time performance, focusing on reducing latency, improving computational efficiency, and ensuring scalability for practical applications. \end{enumerate} 

\noindent The other lines of research outlined, focused on developing new data frameworks, exploring novel architectures, integrating explicit linguistic features, and creating new evaluation metrics, represent crucial long-term endeavors that could each constitute a dissertation in their own right. This proposal takes a strategic and focused approach by pursuing the fifth line of research: \textbf{ optimization and integration of an end-to-end pipeline}. This is a deliberate choice to move from the foundational theoretical work to a tangible applied contribution. By leveraging and fine-tuning state-of-the-art pre-trained models, this dissertation directly addresses the immediate practical challenge of building a functional, real-world system. The resulting low-latency, high-accuracy pipeline will not only serve as a foundational, reproducible tool but also act as a necessary precursor and a real-world testbed upon which future research can build, providing the very platform needed for the innovations that the other lines of research seek to develop.
